% According to ISO, holes should be at a distance of 1.2cm from the border. Hence, an inner border of 2*1.2cm
\usepackage[margin=1.6cm,includefoot,inner=2.4cm,top=3cm]{geometry}

\usepackage[utf8]{inputenc}

\usepackage{helvet}
\usepackage[cm]{sfmath}
\renewcommand{\familydefault}{\sfdefault}

\usepackage{enumitem}

\usepackage{graphicx}

\usepackage{eso-pic}
\usepackage{tikz}

\newcommand{\bgcolor}{black!20}
\newcommand{\indexbgcolor}{black}
\newcommand{\indexfgcolor}{white}
\newcommand{\indexletter}{ }

% Usage: \IfEmpty{\expr}{then}{else}
\newcommand{\IfEmpty}[3]{
	\def\temp{#1}\ifx\temp\empty #2 \else #3 \fi
}

% Model the background
\AddToShipoutPictureBG{ \ifthispageodd { 
	\put(-6,0){
		\begin{tikzpicture}
			\node[minimum height=\paperheight,minimum width = 8.1cm,anchor=west,fill=\bgcolor] at (0,0) {};
		\end{tikzpicture}
	}
	\put(-6,800){
		\begin{tikzpicture}
			\node[minimum width=4cm,fill=\indexbgcolor,text width=3.5cm,align=right,minimum height=.8cm] at (0,0) {\Huge \textcolor{\indexfgcolor}{\textbf \indexletter}};
		\end{tikzpicture}
	}
} {
	\put(375,0) {
		\begin{tikzpicture}
			\node[minimum height=\paperheight,minimum width = 8.1cm,anchor=west,fill=\bgcolor] at (0,0) {};
		\end{tikzpicture}
	}
	\put(475,800){
		\begin{tikzpicture}
			\node[minimum width=4cm,fill=\indexbgcolor,text width=3.5cm,align=left,minimum height=.8cm] at (0,0) {\Huge \textcolor{\indexfgcolor}{\textbf \indexletter}};
		\end{tikzpicture}
	}
} }

\newcounter{recipe}[section]

% Usage: \begin{recipe}{Leckerer Kuchen}{\begin{ingredientcat}...\end{ingredientcat}}{\step{}...\hints{}}
\newcommand{\recipe}[4][]
	{ \begin{minipage}{\textwidth}
		{\noindent\huge \textbf{#2}} \hfill \stepcounter{recipe}\therecipe \\ \rule{\textwidth}{1pt}
		\IfEmpty{#1}{
			\\[\bigskipamount]%
		} {%
			\begin{flushright} \vspace{-\medskipamount} \small Ergibt #1 \end{flushright}%
		}%
		\ifthispageodd {
			\begin{minipage}[t]{.3\textwidth} #3 \end{minipage} \hspace{.03\textwidth}
			\begin{minipage}[t]{.65\textwidth}
				 \noindent {\large \textbf{Anleitung}} \\[-\medskipamount] \rule{\textwidth}{.5pt} 
				\begin{enumerate} #4 \end{enumerate}
			\end{minipage} \\[\bigskipamount]
		} {
			\begin{minipage}[t]{.65\textwidth}
				\noindent {\large \textbf{Anleitung}} \\[-\medskipamount] \rule{\textwidth}{.5pt} 
				\begin{enumerate} #4 \end{enumerate}
			\end{minipage}
			\hspace{.03\textwidth}
			\begin{minipage}[t]{.3\textwidth} #3 \end{minipage} \\[\bigskipamount]
		}
	\end{minipage} }

% Usage: \begin{category}{Teig} \ingredient... \end{category}
\newenvironment{category}[1][]
	{ \IfEmpty{#1}{}{{\large \textbf{#1}}}%
		\begin{itemize}[noitemsep] }
 	{ \end{itemize} }
	
% Usage: \hints{Der Ofen wird angemacht, indem man den Pröngel dreht}
\newcommand{\hints}
	{ \end{enumerate} \noindent {\large \textbf{Hinweise}} \\[-\medskipamount] \rule{\textwidth}{.5pt} \begin{enumerate} \item[]
	}
	
% Usage: \step{Backen}{Backofen öffnen und Zeug reinschieben}
\newcommand{\step}[2][] { \item \IfEmpty{#1}{}{\textbf{#1}} #2 }

% Usage: \ingredient{200}{g}{Zucker} 	
\newcommand{\ingredient}[3]{%
	\item%
	\IfEmpty{#1}{}{$#1$}
	\IfEmpty{#2}{}{#2}
	#3
}